%!TEX root = ./main.tex
%
There is a wide range of service-oriented
architectures (SOAs) dictating design principles for SOC, each one
with its own
idiosyncrasy~\cite{mulesoft:soa,ibm:soa,microsoft:soa,oracle:soa}.
%
We embrace those that hinge on three main concepts: a \emph{service
  provider}, a \emph{service client}, and a \emph{service broker}.
%
The latter handles a \emph{service repository}, a catalogue of service
descriptions searched for in order to discover services required at
runtime.
%
In fact, the service broker is instrumental to the \emph{discovery} of
services according to a contract and of their \emph{binding}, the
composition mechanism that permits to \quo{glue} services together at
runtime as advocated by some SOAs.
% %
% It is worth noticing that in general SOAs do not state how the service
% is discovered or selected from the registry; therefore binding
% mechanisms may require manual intervention.

To support SOC, \SEArch offers a mechanism for populating registries
and composing service-based applications.
%
Registering a service is, in principle, very simple: the service
provider sends the service broker a request for registering a service
attaching a (signed) package containing the contract and the unique
resource identifier (URI) of the provided service.

%%%%%%%%%%%%%%%%%%%%%%%%%%%%%
%%
%%     Una metodología para la certificación de servicios basada en Extended CFSMs
%%
%%%%%%%%%%%%%%%%%%%%%%%%%%%%%
%
%\subsection{A methodology for certification of services based on Extended CFSMs}
% %%% Explica autocertificación de comportamiento de servicios...
%
%%%%%%%%%%%%%%%%%%%%%%%%%%%%%

%%%%%%%%%%%%%%%%%%%%%%%%%%%%%
%%
%%     Certificación de servicios a través de una autoridad certificante
%%
%%%%%%%%%%%%%%%%%%%%%%%%%%%%%
%
%\subsection{Registration of services through a certification authority}
%\Cref{fig:registration} shows how the service registration process could be done if it were to be audited by a certification authority. 
% \begin{figure}[ht]
%  \centering
%  \includegraphics[width=\textwidth]{images/SEArch registration no-ref}
%  \caption{Service registration process}
%  \label{fig:registration}
%\end{figure}
%
%%%%%%%%%%%%%%%%%%%%%%%%%%%%%
 %
 %One possible way of implementing the certification procedure can be a follows. In \SEArch, the provision contract is understood as a Extended CFSM so, if we consider that the functional aspects of values exchanged by the components are expressed in terms of first order formulae, and that those formulae have to be verified to hold under every execution of the implementation, certification can only be achieved up to certain degree of certainty. The rigorous definition of the verification process is not within the scope of the present work; nonetheless, it is possible to sketch it. One can start by synthesising a program skeleton from the Extended CFSM~\cite{lopezpombo:coordination25}, in it, computations take place in the voids generated by the states, those voids reside between communication actions so, it is possible to label the entry point of the voids with first order formulae resulting from the aggregation of conditions in the control flow graph whose arcs correspond to the arrows in the underlying CFSM, then program verification of the code filling the those voids, with respect to the 
%
%%%%%%%%%%%%%%%%%%%%%%%%%%%%%

The execution process of a service-based system in \SEArch is significantly more complex.
%
The computational model behind \SEArch is supported by \emph{Asynchronous Relational Networks}~(ARNs)~\cite{fiadeiro:tcs-503} as the formalisation of the elements of an interface theory for service-oriented software architectures. 
%
  In ARNs there are two types of elements, \emph{processes} and
  \emph{communication channels} both equipped with \emph{ports}.
%
  Ports of processes are called \emph{provides-points},
  while those of communication channels are called
  \emph{requires-points}.
  % : the former constitute the interface through which a service exports its functionality while the latter yield the interface through which an activity expects a certain service to provide a functionality.
  %
  In the operational semantics of ARNs given
  in~\cite{vissani:wadt14-f} actions performed by a component can
  dynamically trigger an automatic and transparent process of
  discovery and binding of a compliant service.
%
The composition of ARNs (i.e., how binding is viewed from a formal perspective) is obtained by \quo{fusing} provides-points and requires-points, subject to a certain compliance check between the contracts associated to them. Under this approach, interoperability is understood at a more abstract level decoupled from any computational aspect.
%
We represent both provision and requirement contracts as \emph{extended Communicating Finite State Machines}~\cite{lopezpombo:coordination25} (extended CFSM for short), an extension of \emph{communicating finite-state machines}~\cite{brand:jacm-30_2} with formal specifications of the functional and non-functional behaviour of service provided / required (depending of the nature of the port); providing a semantic notion of compatibility between provision and
requirement contracts based on a bisimilarity check.

\Cref{fig:execution} depicts the workflow.
 \begin{figure}[t]
  \centering
  \includegraphics[width=0.66\textwidth]{SEArch_execution_no-ref}
  \caption{Service execution procedure in \SEArch}
  \label{fig:execution}
\end{figure}
%
When launched, a component registers its communication channels to
its \emph{middleware}, each of which has its corresponding contract
formalised as a set of extended CFSMs, one for each requires-point.
%
This is required because the middleware has to mediate the
communication with other components.
%
In fact, when the component of a running application, say $C$, tries
to interact with another component over a communication channel, the middleware, say $M$,
captures the attempt and checks whether the communication session for
that communication channel has already been created.
%
If no such session exists, the dynamic reconfiguration process is
triggered as follows:
%
\begin{enumerate}
\item $M$ sends the service broker the contract of the
  communication channel;
\item for each extended CFSM $\mathtt{R}$ in the contract, the
  service broker queries the service repository for candidates;
\item the service repository returns a list of candidates in the form
  $\langle \mathtt{Pr}, u \rangle$, where $\mathtt{Pr}$ is a
  extended CFSM and $u$ is the URI of the service\footnote{\SEArch
	 is parametric in the implementation of the service repository so
	 we assume it is not capable of checking compliance using
	 behavioural contracts. We only rely on its capability of returning
	 a list of candidates, obtained by using potentially more efficient
	 and less precise criteria, for example, an ontology.};
  % the service broker checks whether the provision contract
  % $\mathtt{Pr}$ is bisimilar to the requirement contract
  % $\mathtt{R}$;
  the broker checks whether the provision contract $\mathtt{Pr}$ is
  compatible with the requirement contract $\mathtt{R}$, which is done
  by resorting to bisimilarity check;
\item once the service broker has found services satisfying all the
  requirement contracts, it returns the set of URIs to $M$;
\item $M$ opens a communication with the middleware of each service
  returned by the provider requiring the execution of
  the corresponding service.
\end{enumerate}
%
Then, $M$ sends to or receives from the middleware of the
partner component the actual message; and the execution process
proceeds.
%
Notice that during the execution new requirements might crop up; for
instance, because some brokered services need further
services.\footnote{%
  This implies that a service may have one or more required-points
  associated to communication channels of its own.
	 %
}
% 
This will initiate a new brokerage phase to discover the newly
required services.
% Notice that the execution of the service might also have its own
% requirements putting it as the originator a new dynamic
% reconfiguration.
%
The schematic view discussed before establishes several requisites over
the implementation of middleware and the service broker.
%
We organise the discussion by considering these elements and their
role in the execution architecture:

\paragraph{{\bf The middleware}} provides a private and a public
interface.
%
The former implements functionalities accessible by service clients
and service providers.
%
The public interface implements the capabilities needed for
interacting with service brokers and other middlewares.

The private interface consists of:
\begin{itemize}
\item \texttt{RegisterApp} to register a service and expose it in the
  execution infrastructure. This functionality opens a bidirectional
  (low level) communication channel with the middleware that
  will remain open in order to support the (high level) communication
  with other services.
\item \texttt{RegisterChannel} to register communication channels expressing requirements.
  %
  This functionality  provides the middleware with the relevant information for triggering
  the reconfiguration of the system and managing the communications.
  %
  The functionality can be used by any software artefact running in
  the host, regardless of whether it is a service or client application.
\item \texttt{AppSend} / \texttt{AppRecv} to communicate with partner
  components
\item \texttt{CloseChannel} to close a communication channel.
\end{itemize}
The reader should note the asymmetry between the existence of a
function for explicitly closing a communication session, and the lack
of one for opening it.
%
The reason for this asymmetry resides in that, on the one hand,
transparency in the dynamic reconfiguration of the system is a key
feature of \SEArch but, on the other hand, it is in general not
possible to determine whether a communication session will be used in
the future.

%rpc RegisterChannel(RegisterChannelRequest) returns(RegisterChannelResponse);
%rpc RegisterApp(RegisterAppRequest) returns (stream RegisterAppResponse);
%rpc CloseChannel(CloseChannelRequest) returns (CloseChannelResponse);
%rpc AppSend(AppSendRequest) returns (AppSendResponse);
%rpc AppRecv(AppRecvRequest) returns (AppRecvResponse)

The public interface consists of:
\begin{itemize}
\item \texttt{InitChannel} to accept the initiation of a
  point-to-point (low-level) communication channel.
  %
  This operation allows the broker to initiate the communication
  infrastructure that will connect the service executing behind their
  middleware to the other participants in a communication session
  being setup.
\item \texttt{StartChannel} to receive notifications about point-to-point
  communication channels.
  %
  This operation formally notifies the middleware that the
  brokerage of participants according to a communication channel
  description was successful and the communication session has been
  properly setup.
\item \texttt{MessageExchange} to exchange messages between middlewares.
\end{itemize}

%rpc InitChannel(InitChannelRequest) returns(InitChannelResponse);
%rpc StartChannel(StartChannelRequest) returns(StartChannelResponse);
%rpc MessageExchange (stream MessageExchangeRequest) returns (MessageExchangeResponse);

%\Cref{fig:middleware-design,fig:point-to-point-communication}
%respectively show the structural design of the service middleware and
%the application infrastructure and how buffers are used to provide
%point-to-point communication with external services.
%%
%It is worth noticing that, if the communication channel involve more
%participants, the substructure consisting of Inbox and Outbox buffers,
%and the Sender, will appear replicated.
%\begin{figure}[h]
%\begin{center}
%\subfigure[\label{fig:middleware-design}]{
% \includegraphics[width=0.31\textwidth]{images/SEArch middleware design}
%}
%\subfigure[\label{fig:point-to-point-communication}]{
% \includegraphics[width=0.64\textwidth]{images/SEArch point-to-point communication}
%}
%\caption{Middleware structural design and point-to-point communication.}
%\label{fig:middleware}
%\end{center}
%\end{figure}

\Cref{fig:point-to-point-communication} shows the application infrastructure and how buffers are used to provide
point-to-point communication with external services. Within the infrastructure, it is possible to identify the structural design of the middleware.
 \begin{figure}[ht]
  \centering
  \includegraphics[width=\textwidth]{SEArch_point-to-point_communication}
  \caption{Point-to-point communication between a service client and a service provider.}
  \label{fig:point-to-point-communication}
\end{figure}


\paragraph{{\bf The service broker}}
exposes only two functionalities in a public interface:
\begin{itemize}
\item \texttt{BrokerChannel} to issue requests for brokerage.
  %
  This operation allows a middleware to request for the brokerage of a
  communication channel, and the subsequent creation of a
  communication session to let services interact.
\item \texttt{RegisterProvider} to issue requests for the registration
  of a service provider.
  %
  This operation is the external counterpart of the functionality
  \texttt{RegisterApp} through which the middleware provides
  the service providers the possibility of being offered as services
  available in the execution infrastructure.
\end{itemize}

%rpc BrokerChannel (BrokerChannelRequest) returns (BrokerChannelResponse) {}
%rpc RegisterProvider (RegisterProviderRequest) returns (RegisterProviderResponse) {}

\Cref{fig:seq-diagram-registration} shows the sequence diagram offering a high level view of
the process of registration of a service to the service broker.
\begin{figure}[ht]
  \centering
  \includegraphics[width=0.7\textwidth]{SEArch_registration_seq-diagram}
  \caption{Sequence diagram of the process of registration of a service.}
  \label{fig:seq-diagram-registration}
\end{figure}

\Cref{fig:seq-diagram-brokerage} shows the brokerage process of a
communication channel given the interfaces of the middleware and the
service broker, detailed above in this section.
 \begin{figure}[t!]
  \centering
  \includegraphics[width=\textwidth]{SEArch_brokerage_seq-diagram}
  \caption{Sequence diagram of the process of brokerage of a communication channel.}
  \label{fig:seq-diagram-brokerage}
\end{figure}

The process of brokering a communication channel for building a
session (cf. \cref{fig:seq-diagram-brokerage}) is significantly more
complex.
%
The service client uses the communication channel in the message
\texttt{AppSend} (step \circled{3}).
%
Concurrently, the middleware begins the brokerage
process by sending the contract to the service broker 
(step \circled{4}) and
%
enqueues the message while acknowledging the service client with a
message of type \texttt{AppSendRespond} (step \circled{5}).
%
If the service client has to receive a message (\texttt{AppRecv}), the
middleware captures the attempt triggering the brokerage and going
through the same process for initiating the communication session.
%
In this case, the service client will remain blocked until the
expected message arrives.

The service broker, upon receiving the contract, queries the service
repository for candidates and executes the compliance checks.
%
Each compliance check can be too costly so the service broker
implements a cache for storing precomputed positive results.

After choosing concrete providers for the participants in the
contract, the service broker performs two successive rounds of
message exchanges with the chosen service provider.
%
In the first round, a message of type \texttt{InitChannelRequest} is
sent (step \circled{7}) to tell the middlewares that a
communication session involving its service provider is being
initiated; the message also contains the URIs of all the other
participants in the communication session.
%
At the same time, this message allows the service broker to verify
that the provider is indeed online.
%
Upon receipt of this message, a middleware must accept
incoming messages for this channel and enqueue them for the eventual
reception by the service provider.
%
If all the service providers respond successfully with a message of
type \texttt{InitChannelResponse} (step \circled{8}), then the service
broker performs a second round with a message of type
\texttt{StartChannelRequest} (step \circled{9}), to confirm that the
communication session has been initiated.

After receiving both initialization messages, the middleware
sends the service provider a message of type
\texttt{RegisterAppResponse} (step \circled{{\scriptsize 11}}) containing the UUID
of the new communication session.
%
Then, each service provider can start communicating over this session
according to its contracts.
%
Once a session is initiated, the middlewares establish
unidirectional streams with each other to send messages.
%
In \cref{fig:seq-diagram-brokerage} the middleware of the
service client opens a stream with the other middleware by
sending a message of type \texttt{MessageExchange} (step
{\scriptsize\circled{13}}).
%
After the service provider has received the message (steps
\circled{{\scriptsize 14}}-\circled{{\scriptsize 15}}), it sends a
message (steps \circled{{\scriptsize 16}}-\circled{{\scriptsize 17}})
forcing the middleware of the service provider to establish a stream
in the opposite direction (step {\scriptsize\circled{18}}).

Finally, the service client can close the channel by sending a message
of type \texttt{CloseChannelRequest} (step \circled{{\scriptsize 20}}) to its
middleware, closing the stream used to communicate with the
middleware of the service provider.


%%% Local Variables:
%%% mode: latex
%%% TeX-master: "main"
%%% End:
